\documentclass[a4paper,fleqn]{cas-sc}

% References and mathematics
\usepackage[numbers,sort&compress]{natbib}
\usepackage{amsmath,amssymb}

% Figures
\usepackage{graphicx}
\usepackage[export]{adjustbox}
\usepackage{float}
\usepackage{placeins}
\usepackage{capt-of}

% Tables
\usepackage{booktabs}
\usepackage{array}
\usepackage{tabularx}
\usepackage{multirow}

% Algorithms
\usepackage{algorithm}
\usepackage{algpseudocode}

% Figure paths
\graphicspath{{figures/}{./figures/}{./}}
\DeclareGraphicsExtensions{.pdf,.png,.jpg,.jpeg}

% Table column types
\newcolumntype{Y}{>{\raggedright\arraybackslash}X}
\newcolumntype{Z}[1]{>{\centering\arraybackslash}p{#1}}

% A compact non-floating figure command used throughout the paper.
% Usage:
% \paperfig[trim={...},clip]{0.92\linewidth}{fig.pdf}{Caption.}{fig:label}
\newcommand{\paperfig}[5][]{%
  \begin{center}
  \begin{minipage}{\linewidth}
  \centering
  \includegraphics[
    width=#2,
    max height=0.52\textheight,
    keepaspectratio,
    #1
  ]{#3}%
  \vspace{-0.35em}
  \captionof{figure}{#4}
  \label{#5}
  \end{minipage}
  \vspace{-0.65em}
  \end{center}
}

\let\printorcid\relax

\begin{document}

\setlength{\abovedisplayskip}{0pt plus 1pt minus 0pt}
\setlength{\belowdisplayskip}{5pt plus 1pt minus 1pt}
\setlength{\abovedisplayshortskip}{0pt plus 1pt minus 0pt}
\setlength{\belowdisplayshortskip}{5pt plus 1pt minus 1pt}

\shorttitle{Few-Shot DT Instantiation under Deployment Limits}

\title[mode=title]{Adaptation-Aware Model Selection for Few-Shot Digital Twin Instantiation under Deployment Limits}

\begin{abstract}
As an important modeling tool for managing computing centers, digital twin
(DT) platforms reuse pre-trained DT models from existing centers to reduce the
cost of DT instantiation in large-scale computing networks. However, new or
reconfigured centers may differ substantially from existing centers in resource
configurations and workload characteristics while providing only limited
target-specific data for DT instantiation. Although existing few-shot
adaptation methods can reuse pre-trained models with limited target data, they
typically update model parameters within a predefined architecture. For
heterogeneous target centers, however, the source architecture may be
suboptimal, and the instantiated model must also satisfy per-model
operation-count and parameter-count limits. To address this challenge, we
propose Model Selection and Adaptation for Digital Twin Instantiation
(MSA-DTI), an adaptation-aware model-selection algorithm for few-shot DT
instantiation under deployment limits. MSA-DTI first constructs a candidate
model bank by pre-training models with diverse architectures on historical
workload data, filters out infeasible candidates according to the deployment
limits, and adapts the remaining candidates using limited target samples. To
make model selection more conservative with limited target data, it compares
the best-performing alternative with a predefined reference model after
adaptation and requires a preset relative reduction in validation loss for
replacement. Comparative experiments show that MSA-DTI achieves the lowest
mean squared error (MSE) among the evaluated methods, reducing MSE by 14.60\%
on average across paired cases relative to the pre-training-and-fine-tuning
baseline while satisfying the deployment limits in all main cases.
\end{abstract}

\begin{keywords}
digital twin
\sep few-shot adaptation
\sep model instantiation
\sep architecture selection
\sep model deployment limits
\sep heterogeneous computing centers
\end{keywords}

\maketitle


\section{Introduction}
\label{sec:introduction}

As an important modeling tool for managing computing centers, a shared digital
twin (DT) platform maintains center-specific models for monitoring, prediction,
analysis, and decision making \cite{ref1,ref2,ref3,ref4}. In large-scale computing
networks, pre-trained DT models can be reused and adapted instead of
constructing each target model from scratch, reducing DT-instantiation cost
\cite{ref8,ref11}. Each pre-trained model comprises an architecture and
source-trained parameters. The architecture determines the model structure and
deployment-related complexity, whereas the parameters provide the
initialization for target adaptation. However, a new or reconfigured center
may differ from existing centers in resource configurations,
monitored-variable availability, workload patterns, and data distributions
while providing only limited target-specific data. These differences can make
the original architecture suboptimal for the target center. In addition,
large-scale DT platforms impose per-model resource budgets, represented here
by limits on estimated operation count per inference and parameter count.
Target-center instantiation therefore requires selecting a suitable
architecture and adapting its parameters from limited target data under both
deployment limits. Figure~\ref{fig:scenario} illustrates this setting.

Existing few-shot adaptation methods reuse pre-trained models with limited
target data \cite{ref9,ref10} and have also been applied to DT modeling
\cite{ref11}. However, parameter-adaptation methods typically retain the
original architecture. In contrast, deployment-aware architecture-selection
methods consider model performance together with deployment cost
\cite{ref12,ref13,ref26}. Other approaches jointly consider architecture
decisions and few-shot adaptation \cite{ref30,ref31,ref32}, but typically rely
on target-side architecture or adaptation-structure search. Therefore, the key
challenge is to select and adapt a suitable pre-trained DT model already
available on the platform using limited target data while satisfying per-model
deployment limits.

\paperfig{0.78\linewidth}
{fig1.pdf}
{System setting for few-shot digital twin model instantiation with limited
target data and deployment limits.}
{fig:scenario}

To address this challenge, we propose Model Selection and Adaptation for
Digital Twin Instantiation (MSA-DTI), an adaptation-aware model-selection
algorithm for few-shot DT instantiation under deployment limits. MSA-DTI first
constructs a reusable candidate bank of source-trained models with diverse
architectures. For a target center, it filters candidates that violate the
deployment limits and adapts the remaining feasible candidates under a common
budget. It then compares the adapted candidates using validation performance
and retains a predefined reference model unless the best-performing alternative
achieves a sufficient reduction in validation loss. Thus, model selection is
based on post-adaptation performance without target-side architecture or
adaptation-structure search.

The main contributions are summarized as follows:

\begin{itemize}

\item \textbf{Deployment-constrained few-shot DT instantiation formulation.}
We formulate few-shot DT instantiation as jointly selecting a target
architecture from reusable source-trained models and adapting its parameters
to the target center, subject to hard per-model operation-count and
parameter-count limits.

\item \textbf{Feasibility-first adaptation-aware model selection.}
MSA-DTI filters out infeasible candidates before target adaptation and
compares the remaining candidates only after adapting them using the same
target data and update budget.

\item \textbf{Calibrated reference-model replacement rule.}
We introduce a relative-improvement threshold calibrated on disjoint
development data, ensuring that an adapted alternative replaces the
predefined reference model only after achieving the required relative
reduction in validation loss.

\end{itemize}

The remainder of this paper is organized as follows.
Section~\ref{sec:related} reviews related work.
Section~\ref{sec:problem} defines the system model and problem.
Section~\ref{sec:method} presents MSA-DTI.
Section~\ref{sec:experiments} reports the experiments, and
Section~\ref{sec:conclusion} concludes the paper.

\section{Related Work}
\label{sec:related}

\subsection{Digital Twin Modeling and Platform Deployment}

In communication and computing systems, digital twins represent system states
and services while supporting model updates \cite{ref5,ref6,ref7,ref8}.
At the deployment level, related work considers service orchestration and
computation offloading under resource constraints \cite{ref27,ref28}.
However, DT modeling studies typically assume a predefined model structure,
whereas platform-level research focuses on resource decisions. Neither line of
work directly addresses architecture selection and parameter adaptation for a
target center with limited data under deployment limits.

\subsection{Limited-Data Model Adaptation}

Limited-data methods leverage source or prior knowledge for target-model
adaptation \cite{ref9,ref10,ref24}. Recent approaches improve cross-task and
cross-domain adaptation through context modeling and transferable
representations \cite{ref20,ref21,ref22,ref23}. Limited-data learning has also
been applied to DT construction and workload prediction \cite{ref11,ref25}.
Most parameter-adaptation methods operate with a fixed architecture. For
heterogeneous target centers, a fixed architecture may be unsuitable,
motivating selection among reusable source-trained models.

\subsection{Deployment-Aware Architecture Selection}

Deployment-aware neural architecture search evaluates model performance
together with deployment cost \cite{ref12,ref13,ref26}. Few-shot, zero-shot,
and training-free NAS variants reduce the cost of architecture evaluation
during search \cite{ref14,ref15,ref16}; in Ref.~\cite{ref14}, ``few-shot''
denotes evaluation with several sub-supernets rather than adaptation from a
small target support set.

Among meta-learning-based NAS methods, T-NAS learns transferable architecture
knowledge for new-task adaptation \cite{ref31}. H-Meta-NAS further incorporates
hardware constraints into task-adaptive search \cite{ref30}. Neural
fine-tuning search instead starts from a pre-trained backbone and searches over
adaptation structures, including which components to fine-tune or augment
\cite{ref32}. Together, these methods extend task adaptation beyond parameter
updates to architecture or adaptation-structure decisions.

In contrast, MSA-DTI operates on a finite bank of separately source-trained
reusable models and performs deployment-constrained selection based on
post-adaptation performance, without target-stage search for new architectures
or adaptation structures.

\section{System Model and Problem Formulation}
\label{sec:problem}

\subsection{System Setting and Target Model}

We consider a shared DT platform that maintains center-specific models for
heterogeneous computing centers.

Let $c$ denote a new or reconfigured target center. To enable model reuse
across centers with different monitored-variable availability, all observations
follow a common metric schema, with unavailable variables represented by masks
rather than omitted. Let $\mathcal{X}$ and $\mathcal{Y}$ denote the common
input and output spaces, respectively. Let $\mathcal{A}$ denote the
architecture space. For $a\in\mathcal{A}$, let $\Theta(a)$ denote its
associated parameter space. A model is represented as

\begin{equation}
f_{a,\theta}:\mathcal{X}\rightarrow\mathcal{Y},
\qquad
a\in\mathcal{A},\;
\theta\in\Theta(a),
\label{eq:model-family}
\end{equation}

The DT model instantiated for center $c$ is represented as

\begin{equation}
\mathrm{DT}_{c}
=
\left(
a_c,\theta_c
\right),
\qquad
a_c\in\mathcal{A},\;
\theta_c\in\Theta(a_c),
\label{eq:dt-instance}
\end{equation}

where $a_c$ and $\theta_c$ denote the target architecture and target-specific
parameters, respectively.

The target dataset for center $c$ is represented as

\begin{equation}
\mathcal{D}_{c}^{\mathrm{tar}}
=
\left\{
(x_{c,i},y_{c,i})
\right\}_{i=1}^{K_c},
\label{eq:target-data}
\end{equation}

where $x_{c,i}\in\mathcal{X}$, $y_{c,i}\in\mathcal{Y}$, and $K_c$ denotes
the number of samples available during instantiation. The present problem
concerns determining $(a_c,\theta_c)$ from this limited target data; later
model updates are outside the scope of this study.

\subsection{Reusable Models and Deployment Limits}

Let the finite set of $N$ reusable source-trained models be

\begin{equation}
\mathcal{Q}=\{q_n\}_{n=1}^{N},
\qquad
q_n=(a_n,\theta_n^0),
\label{eq:reusable-models}
\end{equation}

where $a_n\in\mathcal{A}$ denotes the architecture and
$\theta_n^0\in\Theta(a_n)$ denotes its source-trained initialization. The
architectures represented in $\mathcal{Q}$ form the set
$\mathcal{A}_{\mathcal{Q}}=\{a_n\}_{n=1}^{N}\subseteq\mathcal{A}$.

For target center $c$, let $B_c^{F}$ and $B_c^{P}$ denote the operation-count
and parameter-count limits, respectively. For
$a\in\mathcal{A}_{\mathcal{Q}}$, let $F_c(a)$ denote the estimated operation
count of a batch-size-one forward pass under a fixed counting rule, and let
$P_c(a)$ denote the exact number of trainable parameters. In this formulation,
parameter adaptation changes only parameter values while keeping the
architecture fixed; therefore, $F_c(a)$ and $P_c(a)$ remain unchanged.

For center $c$, the deployment-feasible reusable model set is

\begin{equation}
\mathcal{Q}_{c}^{\mathrm{fea}}
=
\left\{
q_n\in\mathcal{Q}
\;\middle|\;
F_c(a_n)\leq B_c^{F},
\;
P_c(a_n)\leq B_c^{P}
\right\}.
\label{eq:feasible-candidate-set}
\end{equation}

Both deployment limits are modeled as hard constraints rather than penalty
terms.

\subsection{DT Model Instantiation Problem}

Let $\mathbb{P}_c$ denote the target-center data distribution over
$\mathcal{X}\times\mathcal{Y}$, and let
$\ell:\mathcal{Y}\times\mathcal{Y}\rightarrow\mathbb{R}_{\geq 0}$ denote the
task loss function. The target risk is defined as

\begin{equation}
R_c(a,\theta)
=
\mathbb{E}_{(x,y)\sim\mathbb{P}_c}
\left[
\ell(f_{a,\theta}(x),y)
\right].
\label{eq:target-risk}
\end{equation}

An ideal target model for center $c$ is defined as a solution to the
constrained optimization problem

\begin{equation}
\begin{aligned}
(a_c^\dagger,\theta_c^\dagger)
\in
\operatorname*{arg\,min}_{a\in\mathcal{A}_{\mathcal{Q}},\,
\theta\in\Theta(a)}
&\quad R_c(a,\theta)
\\
\mathrm{s.t.}\quad
&F_c(a)\leq B_c^F,
\\
&P_c(a)\leq B_c^P.
\end{aligned}
\label{eq:instantiation-problem}
\end{equation}

During instantiation, however, $R_c(a,\theta)$ cannot be evaluated directly
because only the limited target dataset $\mathcal{D}_c^{\mathrm{tar}}$ is
available.

\section{The Proposed MSA-DTI Algorithm}
\label{sec:method}

\subsection{Overview}

Figure~\ref{fig:method-overview} summarizes the three-stage MSA-DTI workflow
for addressing the instantiation problem in Eq.~\eqref{eq:instantiation-problem}:
candidate bank construction, dual-limit candidate filtering, and common-budget
adaptation followed by reference-based selection. The first stage combines
architecture screening and source pretraining to construct a reusable candidate
bank and designates one candidate as the reference before margin calibration
and final target evaluation. For each target center, candidates that violate
the deployment limits are removed before adaptation. The remaining candidates
are adapted using the same target support set and update budget, after which
their post-adaptation performance is compared on a common validation set. A
calibrated relative margin then specifies the relative reduction in validation
loss that the best-performing alternative must achieve to replace the adapted
reference.

\enlargethispage{8pt}
\vspace{-8pt}
\paperfig[trim={60bp 680bp 290bp 90bp},clip]{0.88\linewidth}
{fig2.pdf}
{Overall workflow of MSA-DTI: candidate bank construction, dual-limit
candidate filtering, and common-budget adaptation with reference-based
selection.}
{fig:method-overview}

\subsection{Candidate Bank Construction}

Architecture screening uses a dedicated screening pool that is disjoint from
the data used for margin calibration and final target evaluation. A
reference architecture $a^{\mathrm{ref}}\in\mathcal{A}$ is fixed before
screening and retained regardless of the screening outcome. Each screening
case contains three subsets: support for adaptation, validation for selecting
an alternative architecture, and check for screening evidence. For each
screening case, the reference architecture and all compared alternatives must
satisfy the assigned deployment limits. During screening, the reference and
each feasible alternative are adapted from fixed source-trained screening
initializations on the support subset before validation and check losses are
evaluated.

For each alternative architecture, two evidence counts are accumulated across
the screening cases. The architecture is retained if at least one count reaches
a preset threshold. A \emph{check-oracle win} is assigned to the feasible
alternative architecture with the lowest check-set loss if its loss is lower
than that of the adapted reference. A
\emph{validation-selected positive check win} is assigned to the feasible
alternative architecture selected on the validation subset if its check-set
loss is lower than that of the adapted reference.

Let $\mathcal{D}_{\mathrm{src}}$ denote the combined historical data from the
existing centers. For each retained architecture $a_q$, source pretraining on
$\mathcal{D}_{\mathrm{src}}$ produces an initialization $\theta_q^0$, forming
a candidate $q=(a_q,\theta_q^0)\in\mathcal{Q}$. For a retained architecture
with multiple source-trained initializations, each initialization defines a
distinct candidate. One candidate associated with the reference architecture
$a^{\mathrm{ref}}$ is designated as $q^{\mathrm{ref}}$.
Figure~\ref{fig:source-bank} summarizes the retained architectures and
resulting candidate bank. The candidate bank and reference candidate are fixed
before margin calibration and final target evaluation.

\paperfig[trim={150bp 1340bp 480bp 140bp},clip]{0.92\linewidth}
{fig3.pdf}
{Candidate bank construction and reference candidate selection.}
{fig:source-bank}

\subsection{Dual-Limit Candidate Filtering}

For target center $c$, MSA-DTI computes $\mathcal{Q}_c^{\mathrm{fea}}$ using
Eq.~\eqref{eq:feasible-candidate-set} and restricts target adaptation to this
set. If $q^{\mathrm{ref}}\notin\mathcal{Q}_c^{\mathrm{fea}}$, the procedure
terminates without DT instantiation under the current deployment limits.

\subsection{Common-Budget Adaptation and Reference-Based Selection}

For the target-side procedure, $\mathcal{D}_c^{\mathrm{tar}}$ is partitioned
into a support set and a validation set,
\[
\mathcal{D}_c^{\mathrm{tar}}
=
\mathcal{D}_c^{\mathrm{sup}}\cup\mathcal{D}_c^{\mathrm{val}}.
\]
Support data are used for adaptation, whereas validation data are used for
post-adaptation candidate comparison. Let
$\mathcal{L}(\mathcal{D};a,\theta)$ be the empirical task loss on dataset
$\mathcal{D}$. Each feasible candidate starts from $\theta_q^0$ and is adapted
with a constant learning rate $\eta$:

\begin{equation}
\begin{aligned}
\theta_{c,q}^{(0)}
&=
\theta_q^{0},
\\
\theta_{c,q}^{(t+1)}
&=
\theta_{c,q}^{(t)}
-
\eta
\nabla_{\theta}
\mathcal{L}
\left(
\mathcal{D}_{c}^{\mathrm{sup}};
a_q,
\theta_{c,q}^{(t)}
\right),
\\[-1mm]
&\hspace{38mm}
t=0,\ldots,T-1.
\end{aligned}
\label{eq:target-update}
\end{equation}

Here, $t$ indexes the target-side adaptation steps, $T$ is the common update
budget, and $\eta$ is the learning rate shared by all feasible candidates.
Their post-adaptation validation score is

\begin{equation}
J_c(q)
=
\mathcal{L}
\left(
\mathcal{D}_{c}^{\mathrm{val}};
a_q,
\theta_{c,q}^{(T)}
\right).
\label{eq:validation-loss}
\end{equation}

Because the validation set is small, MSA-DTI uses a relative margin so that a
marginal post-adaptation loss reduction does not by itself trigger replacement
of the fixed reference.

Let
$\mathcal{Q}_{c}^{\mathrm{alt}}
=
\mathcal{Q}_{c}^{\mathrm{fea}}
\setminus
\{q^{\mathrm{ref}}\}$
denote the feasible alternatives. If this set is empty, MSA-DTI sets
$q_c^*=q^{\mathrm{ref}}$ and skips alternative selection. Otherwise, let
$\widetilde q_c
\in
\operatorname*{arg\,min}_{q\in\mathcal{Q}_{c}^{\mathrm{alt}}}
J_c(q)$
denote an alternative with the lowest validation loss.

Let $\tau\in[0,1)$ denote the relative margin. The final candidate is

\begin{equation}
q_c^{*}
=
\begin{cases}
\widetilde q_c,
&
J_c(\widetilde q_c)
\leq
(1-\tau)J_c(q^{\mathrm{ref}}),
\\
q^{\mathrm{ref}},
&
\text{otherwise}.
\end{cases}
\label{eq:reference-based-selection}
\end{equation}

A larger $\tau$ requires a larger validation-loss reduction for an
alternative to be selected over the reference. The margin affects only model
selection and does not change the adapted parameters.

The selected candidate provides the architecture and parameters in
Eq.~\eqref{eq:dt-instance}. Figure~\ref{fig:adaptation-selection} illustrates
the adaptation and selection stages.

\paperfig[trim={140bp 570bp 190bp 120bp},clip]{0.88\linewidth}
{fig5.pdf}
{Common-budget adaptation and reference-based selection.}
{fig:adaptation-selection}

Algorithm~\ref{alg:msa-dti} summarizes the target-side procedure.

\begin{algorithm}[H]
\caption{Target-Side MSA-DTI for One Center}
\label{alg:msa-dti}
\begin{algorithmic}[1]

\Require Candidate bank $\mathcal{Q}$, reference candidate
$q^{\mathrm{ref}}$, support set $\mathcal{D}_{c}^{\mathrm{sup}}$,
validation set $\mathcal{D}_{c}^{\mathrm{val}}$, deployment limits
$B_c^{F}$ and $B_c^{P}$, update budget $T$,
learning rate $\eta$, and relative margin $\tau$

\Ensure Target-specific DT model $\mathrm{DT}_c$, if the reference candidate is deployment-feasible

\State Compute $\mathcal{Q}_{c}^{\mathrm{fea}}$ using
Eq.~\eqref{eq:feasible-candidate-set}

\If{$q^{\mathrm{ref}}\notin\mathcal{Q}_{c}^{\mathrm{fea}}$}
    \State terminate without DT instantiation
\EndIf

\ForAll{$q\in\mathcal{Q}_{c}^{\mathrm{fea}}$}
    \State Adapt $q$ using Eq.~\eqref{eq:target-update}
    \State Compute $J_c(q)$ using Eq.~\eqref{eq:validation-loss}
\EndFor

\State
$\mathcal{Q}_{c}^{\mathrm{alt}}
\gets
\mathcal{Q}_{c}^{\mathrm{fea}}
\setminus
\{q^{\mathrm{ref}}\}$

\If{$\mathcal{Q}_{c}^{\mathrm{alt}}=\emptyset$}
    \State $q_c^{*}\gets q^{\mathrm{ref}}$
\Else
    \State Choose
    $\widetilde q_c
    \in
    \underset{q\in\mathcal{Q}_{c}^{\mathrm{alt}}}
    {\arg\min}\;J_c(q)$
    \State Select $q_c^{*}$ using
    Eq.~\eqref{eq:reference-based-selection}
\EndIf

\State
$\mathrm{DT}_{c}
\gets
\left(
a_{q_c^{*}},
\theta_{c,q_c^{*}}^{(T)}
\right)$

\State \Return $\mathrm{DT}_{c}$

\end{algorithmic}
\end{algorithm}

\subsection{Relative-Margin Calibration}

The margin $\tau$ is calibrated on a margin-calibration pool that is disjoint
from the architecture-screening cases and final target evaluation. Each
calibration case contains support, validation, and check sets. A selection is
harmful when the selected alternative has a higher check-set loss than the
adapted reference. The paired-initialization control uses a second
source-trained initialization of the reference architecture.

Let $\mathcal{T}$ be a preset margin set. For $\tau\in\mathcal{T}$, let
$\widehat R_{\mathrm{harm}}(\tau)$ denote the observed harmful-selection rate,
and let $\widehat G_{\mathrm{ref}}(\tau)$ and
$\widehat G_{\mathrm{pair}}(\tau)$ denote the mean check-loss gains relative
to the adapted reference and paired-initialization control, respectively.
Their lower confidence bounds are $\underline G_{\mathrm{ref}}(\tau)$ and
$\underline G_{\mathrm{pair}}(\tau)$. Given a harmful-selection-rate threshold
$\rho$ and mean-gain thresholds $\gamma_{\mathrm{ref}}$ and
$\gamma_{\mathrm{pair}}$, the eligible margins are

\begin{equation}
\begin{aligned}
\mathcal{T}_{\mathrm{elig}}
=
\bigl\{
\tau\in\mathcal{T}
\;\bigm|\;&
\widehat R_{\mathrm{harm}}(\tau)\leq\rho,
\\
&
\widehat G_{\mathrm{ref}}(\tau)\geq\gamma_{\mathrm{ref}},
\quad
\underline G_{\mathrm{ref}}(\tau)>0,
\\
&
\widehat G_{\mathrm{pair}}(\tau)\geq\gamma_{\mathrm{pair}},
\quad
\underline G_{\mathrm{pair}}(\tau)>0
\bigr\}.
\end{aligned}
\label{eq:margin-calibration}
\end{equation}

When $\mathcal{T}_{\mathrm{elig}}$ is nonempty, MSA-DTI sets
\[
\tau^*=\min\mathcal{T}_{\mathrm{elig}}.
\]
The margin grid, threshold values, and resulting $\tau^*$ are reported in
Section~\ref{sec:experiments}. The calibrated margin is fixed before final
target evaluation.

If $\mathcal{T}_{\mathrm{elig}}=\varnothing$, MSA-DTI retains the adapted
reference. Because the statistics in Eq.~\eqref{eq:margin-calibration} are
estimated from the calibration pool, the calibrated margin depends on the
calibration domain.

\subsection{Analytical Properties}

\textbf{Deployment feasibility.}
Any DT model instantiated by MSA-DTI is selected from
$\mathcal{Q}_{c}^{\mathrm{fea}}$:

\[
F_c(a_{q_c^{*}})\leq B_c^{F},
\qquad
P_c(a_{q_c^{*}})\leq B_c^{P}.
\]

\textbf{Required validation improvement.}
When an alternative is selected and $J_c(q^{\mathrm{ref}})>0$,
Eq.~\eqref{eq:reference-based-selection} implies

\[
\frac{
J_c(q^{\mathrm{ref}})
-
J_c(q_c^{*})
}{
J_c(q^{\mathrm{ref}})
}
\geq
\tau.
\]

The target-stage computational complexity is
\[
O\left(
|\mathcal{Q}|+
\sum_{q\in\mathcal{Q}_{c}^{\mathrm{fea}}}
\left(TC_q^{\mathrm{grad}}+C_q^{\mathrm{val}}\right)
\right),
\]
where $C_q^{\mathrm{grad}}$ and $C_q^{\mathrm{val}}$ denote the costs of one
target update and validation evaluation for candidate $q$, respectively, while
the $|\mathcal{Q}|$ term accounts for deployment filtering.

\section{Experimental Evaluation}
\label{sec:experiments}

We evaluate prediction accuracy, selection reliability, deployment
feasibility, target-side cost, and robustness across target conditions.
Additional component and transfer analyses are reported.

\subsection{Experimental Setup}
\label{subsec:experimental-setup}

\subsubsection{Benchmark and Data Protocol}

We use a controlled synthetic benchmark for heterogeneous multi-center
workloads. Center types A--C are evaluation strata with progressively sparser
and noisier monitoring and distinct workload-dynamics ranges; deployment-limit
tiers are separate platform-side settings.

The main study uses 20 source centers and 20 held-out target centers.
Prediction horizons $H\in\{1,4\}$ and support sizes $K\in\{10,20\}$ give
four cases per target center and 80 held-out cases in total.

Each case is split chronologically into support, validation, check, and test
sets, with gaps between adjacent splits. Check sets are used for architecture
screening, margin calibration, and the separate diagnostic analyses; held-out
target test data are used only after model selection is fixed. The source,
architecture-screening, margin-calibration, diagnostic, and held-out target
pools use disjoint center identities.

\subsubsection{Candidate Bank, Reference, and Calibration}
\label{subsubsec:bank-calibration}

The initial architecture space contains 66 configurations spanning multilayer
perceptron (MLP), temporal convolutional network (TCN), and gated recurrent
unit (GRU) architectures. A57, a one-layer GRU with 32 hidden units, is fixed
as the reference before screening.

An alternative configuration is retained if either of the two screening
evidence counts defined in Section~4.2 is at least two. The screening retains six configurations corresponding to five executable
architectures, and the frozen source-trained records form a seven-candidate
bank.

Applying Eq.~\eqref{eq:margin-calibration} to the preset margin grid on the
main calibration pool yields the smallest eligible margin $\tau=0.10$, which
is frozen before held-out evaluation. Table~\ref{tab:configuration}
summarizes the experimental configuration.

\begin{center}
\begin{minipage}{\linewidth}
\captionof{table}{Main experimental configuration.}
\label{tab:configuration}
\centering
\scriptsize
\renewcommand{\arraystretch}{1.05}
\setlength{\tabcolsep}{2.5pt}
\begin{tabularx}{\linewidth}{Z{0.18\linewidth} Y Z{0.43\linewidth}}
\toprule
Category & Item & Setting \\
\midrule
Data
& Centers; observations; input length
& Source / held-out: $20/20$; 3000 per center; $L=96$ \\
& Target cases and data
& 80 cases; $H\in\{1,4\}$; $K\in\{10,20\}$;
validation / check / test: $80/60/200$ \\
\midrule
Bank
& Initial / retained / executable / candidates
& $66/6/5/7$ \\
& Reference and screening
& A57: one-layer GRU-32, fixed before screening; alternative retained if
either screening evidence count in Section~4.2 is $\geq 2$ \\
& Source-model records
& Reference pooled-source initialization; five common alternative
initializations; one additional A57 paired control \\
\midrule
Training
& Source models
& Adam / MSE / 50 epochs; learning rate $10^{-3}$; batch size 64 \\
& Target adaptation
& SGD / MSE / 50 updates; learning rate 0.01; gradient-norm limit 1.0;
complete support set \\
\midrule
Selection
& Grid and eligibility
& $\{0.05,0.075,0.10,0.125,0.15,0.20\}$;
harmful-selection rate $\leq5\%$; both mean gains $\geq3\%$;
both 95\% lower confidence bounds $>0$ \\
& Frozen margin
& $\tau=0.10$ \\
\midrule
Limits
& Operation-count limits
& Tight / medium / loose:
$1.5{\times}10^6/5{\times}10^6/2{\times}10^7$ \\
& Parameter-count limits
& Tight / medium / loose:
$3{\times}10^4/10^5/5{\times}10^5$ \\
\midrule
Statistics
& Confidence intervals
& Center-cluster bootstrap; 4000 repetitions \\
\bottomrule
\end{tabularx}
\end{minipage}
\end{center}


\subsubsection{Task, Baselines, and Metrics}

We use workload prediction as the evaluation task because it supports
proactive resource management in computing systems
\cite{ref17,ref18,ref19}. Mean squared error (MSE) is the task loss. We also
report mean absolute error (MAE), Worst-10\% error, conditional value at risk
at the 90th percentile (CVaR90), target-side time, and limit satisfaction.

For each case, Worst-10\% error is the mean of the largest 10\% of window-level
MSE values on the test set. CVaR90 is the mean case-level MSE over the 10\% of
target cases with the largest MSE.

For held-out reporting, a selected alternative is labeled harmful when its
test MSE is higher than that of the adapted reference. Relative gains are
computed for each paired case and then averaged. The 95\% confidence intervals
use a center-cluster bootstrap with 4000 repetitions, resampling target centers
while keeping their four cases together.

We compare MSA-DTI with pre-training and fine-tuning (PT+FT), MeDeT-based
adaptation, random initialization, few-shot NAS, zero-shot NAS, zero-shot NAS
with fine-tuning, and task-matched H-Meta-NAS-based adaptation
\cite{ref30}. PT+FT is the main controlled baseline because it shares the
source-trained reference and target adaptation procedure with MSA-DTI.

As the closest evaluated architecture-adaptation baseline, H-Meta-NAS-based
adaptation uses source-learned meta-initializations and validation-only target
search after deployment filtering. Its meta-initializations and search settings
are fixed without held-out test data, and the main comparison uses a
budget-matched 50-update SGD/MSE setting. All methods use the same target
cases, chronological splits, deployment limits, and architecture definitions.

\subsubsection{Runtime Protocol}

All methods are implemented in Python and PyTorch and run on an NVIDIA RTX
3060 Laptop GPU. Target-side time includes model creation, initialization
loading, target-side search when applicable, target adaptation, validation
evaluation, and final selection, but excludes offline preparation and test
evaluation. Each reported time is the mean over five runs.

\subsection{Overall Performance and Selection Reliability}
\label{subsec:overall}

Figure~\ref{fig:overall-performance} compares predictive performance on the
held-out target cases. All selected models satisfy both deployment limits.

\paperfig{0.98\linewidth}
{fig_overall_performance_ours.pdf}
{Predictive performance on the held-out target cases: (a) MAE, (b) MSE, (c) Worst-10\% error, and (d) CVaR90 (lower is better).}
{fig:overall-performance}

MSA-DTI gives the lowest values for all four error measures. Compared with
PT+FT, the strongest baseline, MSA-DTI reduces MAE by 8.55\%, MSE by 14.60\%,
and Worst-10\% error by 13.76\% on average across paired cases. The 95\%
confidence interval for the paired MSE reduction is $[9.41\%,20.14\%]$.
Paired MSE comparisons with the remaining baselines are summarized in
Supplementary Section~S4.2; each 95\% confidence interval lies above zero.

Figure~\ref{fig:paired-instantiation} compares MSA-DTI with PT+FT at the case
level.

\paperfig{0.70\linewidth}
{fig6.pdf}
{Held-out paired comparison of MSA-DTI and PT+FT for (a) test MSE and (b) test Worst-10\% error; beneficial and harmful alternative labels are defined by test MSE, and points below the diagonal favor MSA-DTI.}
{fig:paired-instantiation}

MSA-DTI retains the reference in 33 cases and selects an alternative in 47
cases. Of the 47 alternative selections, 44 have lower test MSE than the
adapted reference and three are harmful, corresponding to 3.75\% of all
held-out cases. The 95\% confidence interval for the all-case
harmful-selection rate is $[0.00\%,7.50\%]$, so it does not establish a
population harmful-selection rate below the 5\% calibration criterion.

Because the native NAS baselines use different target optimizers and losses,
we additionally evaluate an optimizer-matched diagnostic in which the included
methods use SGD, MSE, and 50 target updates per adapted model. In this
diagnostic, MSA-DTI achieves the lowest MSE, Worst-10\% error, and CVaR90 among
the included methods. Table~\ref{tab:optimizer-matched} reports these results.

\begin{center}
\begin{minipage}{\linewidth}
\captionof{table}{Optimizer-matched diagnostic comparison.}
\label{tab:optimizer-matched}
\centering
\scriptsize
\renewcommand{\arraystretch}{1.07}
\setlength{\tabcolsep}{2.8pt}
\begin{tabularx}{\linewidth}
{Y Z{0.17\linewidth} Z{0.20\linewidth} Z{0.18\linewidth} Z{0.18\linewidth}}
\toprule
Method
& MSE$\downarrow$
& Worst-10\%$\downarrow$
& CVaR90$\downarrow$
& Adapted models per case \\
\midrule
PT+FT & 0.004497 & 0.017144 & 0.011878 & 1.00 \\
Few-shot NAS & 0.006715 & 0.025246 & 0.019653 & 12.00 \\
Zero-shot NAS + fine-tuning & 0.007402 & 0.028099 & 0.020402 & 12.00 \\
\textbf{MSA-DTI} & \textbf{0.003809} & \textbf{0.015094} &
\textbf{0.008696} & 6.10 \\
\bottomrule
\end{tabularx}
\vspace{2pt}
\par\footnotesize
All methods use the same independent diagnostic cases, hard deployment
limits, SGD optimizer, MSE loss, and 50 target updates per adapted model. The
source initializations and candidate-selection rules remain method-specific.
\end{minipage}
\end{center}

\subsection{Robustness and Sensitivity Analysis}
\label{subsec:heterogeneity}

Figure~\ref{fig:heterogeneity-robustness} reports paired MSE reductions
relative to PT+FT. Positive values favor MSA-DTI.

\paperfig{0.88\linewidth}
{fig7.pdf}
{Paired MSE reduction relative to PT+FT across (a) horizon/support settings and (b) target-center types; error bars show 95\% confidence intervals.}
{fig:heterogeneity-robustness}

The reduction ranges from 11.57\% to 18.03\% across the four horizon and
support-size settings and from 13.03\% to 16.80\% across the three center
types. All seven confidence intervals remain above zero.

On a separate robustness pool, the fixed candidate architectures retrained
under three source-training seeds yield paired MSE reductions ranging from
10.52\% to 13.21\%. All three 95\% confidence intervals remain above zero.

We additionally apply the full 66-configuration screening procedure to two
new 20-center pools, giving three screening pools in total when the original
pool is included. Each screening outcome is then evaluated on a fourth,
untouched 20-center pool. The retained-list sizes are 6, 0, and 4; the two new
lists have Jaccard similarities of 0 and 0.429 to the frozen six-configuration
list. On the independent diagnostic pool, validation-based selection from the
frozen list gives a 0.30\% mean paired check-MSE reduction, with a 95\%
confidence interval of $[-0.65\%,1.56\%]$ and a 2.50\% harmful-selection
rate.

These results show that the preset candidate-discovery rule is
sample-sensitive. The main experiment uses the prespecified bank frozen before
held-out evaluation. In the rescreening analysis, A57 remains the reference
when no alternative configuration is retained.

\subsection{Deployment Feasibility and Cost}
\label{subsec:deployment}

Feasibility in this subsection is assessed using the operation-count and
parameter-count limits defined in Section~3.2. Figure~\ref{fig:budget-architecture}
shows how the three deployment-limit settings affect candidate feasibility and
final model selection.

\paperfig{0.98\linewidth}
{fig8.pdf}
{Effects of the three deployment-limit settings on
(a) the number of feasible candidates and
(b) the selected model configurations.}
{fig:budget-architecture}

Under the tight setting, four of the seven candidates in the frozen bank are
deployment-feasible, whereas all seven are feasible under the medium and loose
settings.

For the frozen main bank under the tight setting, a single-limit analysis
shows that the reduction is caused by the parameter-count limit: all seven
candidates satisfy the operation-count limit, while the parameter-count limit
excludes three MLP candidates.

An expanded-bank stress analysis shows that the two limits exclude different
candidates for both $H=1$ and $H=4$: each single-limit feasible set retains
seven of eight candidates, whereas their joint feasible set retains six.

In the target-side stress evaluation on the same 80 held-out cases, all
selected models satisfy both deployment limits. Mean paired MSE reductions
relative to PT+FT remain positive under tight, medium, and loose limits:
11.35\%, 15.30\%, and 18.22\%, respectively.

Table~\ref{tab:target-cost} compares target-side cost and selected-model
complexity for PT+FT, the target-adaptation architecture-search baselines, and
MSA-DTI.

\begin{center}
\begin{minipage}{\linewidth}
\captionof{table}{
Target-side instantiation time, parameter count, and estimated operation
count. Runtime is reported as mean $\pm$ standard deviation across five GPU-synchronized repeats.
}
\label{tab:runtime_complexity}
\centering
\scriptsize
\renewcommand{\arraystretch}{1.08}
\setlength{\tabcolsep}{3.0pt}
\begin{tabular*}{\linewidth}{@{\extracolsep{\fill}}lcccc@{}}
\toprule
Method
& Time (s)$\downarrow$
& Adapted models
& \shortstack{Parameter\\count$\downarrow$}
& \shortstack{Estimated operation\\count$\downarrow$} \\
\midrule
PT+FT & 0.204 $\pm$ 0.004 & 1.00 & 5,746 & 1,050,784 \\
MeDeT-based adaptation & 0.204 $\pm$ 0.004 & 1.00 & 5,746 & 1,050,784 \\
Random initialization & 0.198 $\pm$ 0.009 & 1.00 & 5,746 & 1,050,784 \\
Few-shot NAS & 13.471 $\pm$ 0.173 & 12.00 & 30,995 & 1,468,894 \\
Zero-shot NAS & 0.972 $\pm$ 0.010 & 0.00 & 86,485 & 481,058 \\
Zero-shot NAS + fine-tuning & 17.601 $\pm$ 0.150 & 12.00 & 62,076 & 1,571,621 \\
\textbf{Proposed method} & 5.676 $\pm$ 0.059 & 6.25 & 14,892 & 751,433 \\
\bottomrule
\end{tabular*}
\vspace{2pt}
\footnotesize
Time covers target-side instantiation only. It excludes data construction,
check and test evaluation, and source training.
\end{minipage}
\end{center}


MSA-DTI is slower than PT+FT but $2.37\times$ faster than few-shot NAS and
$3.10\times$ faster than zero-shot NAS with fine-tuning. The reported
target-side times are measured for one instantiation case on one GPU;
concurrent-case scaling and platform-level scheduling are outside the present
evaluation.

Hardware profiling evaluates every executable architecture on an Intel Core
i7-12700H CPU and an NVIDIA RTX 3060 Laptop GPU. Median batch-size-one
inference latency ranges from 0.093 to 1.980 ms on the CPU and from 0.133 to
4.170 ms on the GPU. Serialized model size ranges from 10.3 to 315.9 KiB.
Operation count is positively associated with median latency, whereas
parameter count preserves the serialized-size ranking in all measured
settings. Peak GPU allocation is not monotonic in parameter count. These
measurements support the two counts as distinct model-level complexity
indicators; their hardware-level implications remain host-dependent.

\subsection{Component Analysis}
\label{subsec:component}

Table~\ref{tab:component} reports component variants on an independent
diagnostic pool. The reference-only and two-reference-initialization variants
do not reproduce the cross-architecture gain. Without the relative margin,
diagnostic MSE is nearly unchanged, but alternatives are selected more often
and the harmful-selection rate increases from 15.00\% to 17.50\%. The
corresponding confidence intervals overlap, so this increase is treated as an
observed rather than statistically established difference. Removing deployment
filtering slightly lowers MSE but reduces limit satisfaction from 100\% to
96.25\%.

Overall, the main held-out result evaluates the complete MSA-DTI pipeline,
while the diagnostic controls distinguish the roles of individual design
choices. The cross-architecture gain is not reproduced by the
same-architecture controls and remains sensitive to source training. The
relative margin mainly affects replacement conservativeness, whereas
deployment filtering enforces limit satisfaction.

\begin{center}
\begin{minipage}{\linewidth}
\captionof{table}{Component variants on the independent diagnostic pool.}
\label{tab:component}
\centering
\scriptsize
\renewcommand{\arraystretch}{1.08}
\setlength{\tabcolsep}{2.2pt}
\begin{tabularx}{\linewidth}
{Y Z{0.12\linewidth} Z{0.14\linewidth} Z{0.17\linewidth} Z{0.16\linewidth} Z{0.16\linewidth}}
\toprule
Variant
& MSE$\downarrow$
& Worst-10\%$\downarrow$
& Paired MSE reduction (\%)$\uparrow$
& Harmful-selection rate (\%)$\downarrow$
& Limit satisfaction (\%)$\uparrow$ \\
\midrule
\textbf{MSA-DTI} & 0.004063 & 0.015764 & 15.79 & 15.00 & 100.00 \\
Earlier source training & 0.004876 & 0.019044 & 1.75 & 2.50 & 100.00 \\
Reference only & 0.005051 & 0.019512 & 0.00 & 0.00 & 100.00 \\
Two reference initializations & 0.004917 & 0.019136 & 1.29 & 1.25 & 100.00 \\
Without relative margin & 0.004055 & 0.015690 & 16.16 & 17.50 & 100.00 \\
Without deployment filtering & 0.004037 & 0.015635 & 16.13 & 13.75 & 96.25 \\
\bottomrule
\end{tabularx}
\vspace{2pt}
\par\footnotesize
Paired MSE reduction is computed relative to the adapted reference on the same
diagnostic pool; harmful-selection rate is computed over all diagnostic
cases. The earlier source-training row changes the full source-training
procedure for the alternatives, so it measures candidate-bank sensitivity
rather than a single-component ablation.
\end{minipage}
\end{center}

\subsection{External Transfer Evaluation}
\label{subsec:external}

We also evaluate MSA-DTI on Alibaba Cluster Trace v2018 \cite{ref29}. The
Alibaba data are divided into disjoint source, calibration, and held-out
machine pools; source models are retrained on the Alibaba source pool while
the method configuration remains fixed. The workload observations come from
the production trace, while the deployment-limit tiers use the same controlled
model-complexity settings as in the main study.

When the main-study margin $\tau=0.10$, fixed before the Alibaba evaluation,
is transferred without recalibration, the mean paired MSE reduction is
4.20\%, with a 95\% confidence interval of $[2.03\%,6.70\%]$. The
harmful-selection rate is 6.25\%, exceeding the original 5\% calibration
criterion. No margin in the preset grid satisfies the calibration criteria on
the Alibaba calibration pool.

Under the calibration rule, the adapted reference is therefore retained. The
fixed $\tau=0.10$ result is reported separately to assess direct transfer of
the main-study margin.

\section{Conclusion}
\label{sec:conclusion}

This paper presented MSA-DTI for few-shot DT model instantiation under
per-model operation-count and parameter-count limits. MSA-DTI filters
infeasible candidates from the source-trained bank, adapts the remaining
candidates under a common update budget, and applies a calibrated
reference-model replacement rule based on post-adaptation validation
performance. In the main held-out evaluation, MSA-DTI reduces MSE by 14.60\%
on average across paired cases relative to PT+FT while satisfying both
deployment limits in all cases. The confidence interval for the held-out
harmful-selection rate spans the 5\% calibration criterion, and the
candidate-discovery rule is sample-sensitive. On Alibaba, direct transfer of
the fixed main-study margin yields a positive mean paired MSE reduction, but
no margin in the preset grid satisfies the calibration criteria. The current
evaluation is limited to a controlled synthetic benchmark and one external
workload trace. Future work includes investigating more stable
candidate-discovery procedures, uncertainty-aware selection, and deployment
limits derived from measurements on additional platforms.

\section*{CRediT Authorship Contribution Statement}

\textit{The final author contribution statement will be added after the
author list is confirmed.}

\section*{Funding}

\textit{Funding information will be added after confirmation.}

\section*{Declaration of Competing Interest}

The authors declare that they have no known competing financial
interests or personal relationships that could have appeared to
influence the work reported in this paper.

\section*{Data Availability}

The source code, frozen configurations, processed results, plot-ready data,
and reproduction scripts underlying the reported results are available in
GitHub Release v1.2.7:
\url{https://github.com/w924871681/paper-code-DT-Twin/releases/tag/v1.2.7}.
The original Alibaba Cluster Trace v2018 is not redistributed and remains
available from its official repository:
\url{https://github.com/alibaba/clusterdata}. The released documentation
specifies preprocessing, disjoint data splits, margin handling, and
held-out evaluation.

\section*{Declaration of Generative AI and AI-Assisted Technologies
in the Manuscript Preparation Process}

During the preparation of this work, the authors used OpenAI ChatGPT to
assist with language refinement, manuscript organization, and readability
improvement. After using this tool, the authors reviewed and edited the
content as needed and take full responsibility for the content of the
published article.

\begin{thebibliography}{00}
\bibitem{ref1} L. U. Khan, Z. Han, W. Saad, E. Hossain, M. Guizani, and C. S. Hong, ``Digital twin of wireless systems: Overview, taxonomy, challenges, and opportunities,'' \emph{IEEE Communications Surveys \& Tutorials}, vol. 24, no. 4, pp. 2230--2254, Fourth Quarter 2022, doi: 10.1109/COMST.2022.3198273.

\bibitem{ref2} S. Mihai, M. Yaqoob, D. V. Hung, W. Davis, P. Towakel, M. Raza, M. Karamanoglu, B. Barn, D. Shetve, R. V. Prasad, H. Venkataraman, R. Trestian, and H. X. Nguyen, ``Digital twins: A survey on enabling technologies, challenges, trends and future prospects,'' \emph{IEEE Communications Surveys \& Tutorials}, vol. 24, no. 4, pp. 2255--2291, Fourth Quarter 2022, doi: 10.1109/COMST.2022.3208773.

\bibitem{ref3} H. Xu, J. Wu, Q. Pan, X. Guan, and M. Guizani, ``A survey on digital twin for Industrial Internet of Things: Applications, technologies and tools,'' \emph{IEEE Communications Surveys \& Tutorials}, vol. 25, no. 4, pp. 2569--2598, Fourth Quarter 2023, doi: 10.1109/COMST.2023.3297395.

\bibitem{ref8} J. Li, S. Guo, W. Liang, J. Wang, Q. Chen, Y. Zeng, B. Ye, and X. Jia, ``Digital twin-enabled service provisioning in edge computing via continual learning,'' \emph{IEEE Transactions on Mobile Computing}, vol. 23, no. 6, pp. 7335--7350, 2024, doi: 10.1109/TMC.2023.3332668.

\bibitem{ref11} H. Sartaj, S. Ali, and J. M. Gj{\o}by, ``MeDeT: Medical device digital twins creation with few-shot meta-learning,'' \emph{ACM Transactions on Software Engineering and Methodology}, vol. 34, no. 6, Art. no. 158, pp. 1--36, 2025, doi: 10.1145/3708534.

\bibitem{ref9} A. Vettoruzzo, M.-R. Bouguelia, J. Vanschoren, T. R{\"o}gnvaldsson, and K. C. Santosh, ``Advances and challenges in meta-learning: A technical review,'' \emph{IEEE Transactions on Pattern Analysis and Machine Intelligence}, vol. 46, no. 7, pp. 4763--4779, 2024, doi: 10.1109/TPAMI.2024.3357847.

\bibitem{ref10} H. Gharoun, F. Momenifar, F. Chen, and A. H. Gandomi, ``Meta-learning approaches for few-shot learning: A survey of recent advances,'' \emph{ACM Computing Surveys}, vol. 56, no. 12, Art. no. 294, pp. 1--41, 2024, doi: 10.1145/3659943.

\bibitem{ref12} M. Tan, B. Chen, R. Pang, V. Vasudevan, M. Sandler, A. Howard, and Q. V. Le, ``MnasNet: Platform-aware neural architecture search for mobile,'' in \emph{Proc. IEEE/CVF Conference on Computer Vision and Pattern Recognition (CVPR)}, 2019, pp. 2820--2828, doi: 10.1109/CVPR.2019.00293.

\bibitem{ref13} B. Wu, X. Dai, P. Zhang, Y. Wang, F. Sun, Y. Wu, Y. Tian, P. Vajda, Y. Jia, and K. Keutzer, ``FBNet: Hardware-aware efficient ConvNet design via differentiable neural architecture search,'' in \emph{Proc. IEEE/CVF Conference on Computer Vision and Pattern Recognition (CVPR)}, 2019, pp. 10734--10742, doi: 10.1109/CVPR.2019.01099.

\bibitem{ref26} Y. Zhao, L. Wang, and T. Guo, ``Multi-objective neural architecture search by learning search space partitions,'' \emph{Journal of Machine Learning Research}, vol. 25, no. 177, pp. 1--41, 2024.

\bibitem{ref30}
Y. Zhao, X. Gao, I. Shumailov, N. Fusi, and R. Mullins,
``Rapid model architecture adaption for meta-learning,''
in \emph{Advances in Neural Information Processing Systems},
vol. 35, 2022,
doi: 10.52202/068431-1360.

\bibitem{ref31}
D. Lian, Y. Zheng, Y. Xu, Y. Lu, L. Lin, P. Zhao, J. Huang, and S. Gao,
``Towards fast adaptation of neural architectures with meta learning,''
in \emph{Proc. International Conference on Learning Representations
(ICLR)}, 2020.

\bibitem{ref32}
P. Eustratiadis, L. Dudziak, D. Li, and T. Hospedales,
``Neural fine-tuning search for few-shot learning,''
in \emph{Proc. International Conference on Learning Representations
(ICLR)}, 2024.


\bibitem{ref4} F. Tao, H. Zhang, A. Liu, and A. Y. C. Nee, ``Digital twin in industry: State-of-the-art,'' \emph{IEEE Transactions on Industrial Informatics}, vol. 15, no. 4, pp. 2405--2415, 2019, doi: 10.1109/TII.2018.2873186.

\bibitem{ref5} B. Qin, H. Pan, Y. Dai, X. Si, X. Huang, C. Yuen, and Y. Zhang, ``Machine and deep learning for digital twin networks: A survey,'' \emph{IEEE Internet of Things Journal}, vol. 11, no. 21, pp. 34694--34716, Nov. 2024, doi: 10.1109/JIOT.2024.3416733.

\bibitem{ref6} A. Hakiri, A. Gokhale, S. Ben Yahia, and N. Mellouli, ``A comprehensive survey on digital twin for future networks and emerging Internet of Things industry,'' \emph{Computer Networks}, vol. 244, Art. no. 110350, 2024, doi: 10.1016/j.comnet.2024.110350.

\bibitem{ref7} Y. Wu, K. Zhang, and Y. Zhang, ``Digital twin networks: A survey,'' \emph{IEEE Internet of Things Journal}, vol. 8, no. 18, pp. 13789--13804, Sept. 2021, doi: 10.1109/JIOT.2021.3079510.

\bibitem{ref27}
M. Emami Khansari and S. Sharifian,
``A deep reinforcement learning approach towards distributed Function as a
Service (FaaS) based edge application orchestration in cloud-edge
continuum,''
\emph{Journal of Network and Computer Applications}, vol. 233,
Art. no. 104042, 2025,
doi: 10.1016/j.jnca.2024.104042.

\bibitem{ref28} D. Hortelano, I. de Miguel, R. J. Dur{\'a}n Barroso,
J. C. Aguado, N. Merayo, L. Ruiz, A. Asensio, X. Masip-Bruin,
P. Fern{\'a}ndez, R. M. Lorenzo, and E. J. Abril,
``A comprehensive survey on reinforcement-learning-based computation
offloading techniques in Edge Computing Systems,''
\emph{Journal of Network and Computer Applications}, vol. 216,
Art. no. 103669, 2023, doi: 10.1016/j.jnca.2023.103669.

\bibitem{ref24} J. Son, S. Lee, and G. Kim, ``When meta-learning meets online and continual learning: A survey,'' \emph{IEEE Transactions on Pattern Analysis and Machine Intelligence}, vol. 47, no. 1, pp. 413--432, 2025, doi: 10.1109/TPAMI.2024.3463709.

\bibitem{ref20} L. M. Zintgraf, K. Shiarlis, V. Kurin, K. Hofmann, and S. Whiteson, ``Fast context adaptation via meta-learning,'' in \emph{Proc. 36th Int. Conf. Mach. Learn. (ICML)}, vol. 97 of \emph{Proc. Mach. Learn. Res.}, pp. 7693--7702, 2019.

\bibitem{ref21} Y. Zhao, T. Zhang, J. Li, and Y. Tian, ``Dual adaptive representation alignment for cross-domain few-shot learning,'' \emph{IEEE Transactions on Pattern Analysis and Machine Intelligence}, vol. 45, no. 10, pp. 11720--11732, 2023, doi: 10.1109/TPAMI.2023.3272697.

\bibitem{ref22} J. Wu, X. Liu, X. Yin, T. Zhang, and Y. Zhang, ``Task-adaptive prompted transformer for cross-domain few-shot learning,'' in \emph{Proc. AAAI Conf. Artif. Intell.}, vol. 38, no. 6, pp. 6012--6020, 2024, doi: 10.1609/aaai.v38i6.28416.

\bibitem{ref23} W. Wang, L. Duan, Y. Wang, J. Fan, and Z. Zhang, ``MMT: Cross domain few-shot learning via meta-memory transfer,'' \emph{IEEE Transactions on Pattern Analysis and Machine Intelligence}, vol. 45, no. 12, pp. 15018--15035, Dec. 2023, doi: 10.1109/TPAMI.2023.3306352.

\bibitem{ref25} Z. Zuo, Y. Huang, Z. Li, Y. Jiang, and C. Liu, ``Mixed contrastive transfer learning for few-shot workload prediction in the cloud,'' \emph{Computing}, vol. 107, no. 1, Art. no. 5, 2025, doi: 10.1007/s00607-024-01366-y.

\bibitem{ref14} Y. Zhao, L. Wang, Y. Tian, R. Fonseca, and T. Guo, ``Few-shot neural architecture search,'' in \emph{Proc. 38th Int. Conf. Mach. Learn. (ICML)}, vol. 139 of \emph{Proc. Mach. Learn. Res.}, pp. 12707--12718, 2021.

\bibitem{ref15} G. Li, D. Hoang, K. Bhardwaj, M. Lin, Z. Wang, and R. Marculescu, ``Zero-shot neural architecture search: Challenges, solutions, and opportunities,'' \emph{IEEE Transactions on Pattern Analysis and Machine Intelligence}, vol. 46, no. 12, pp. 7618--7635, Dec. 2024, doi: 10.1109/TPAMI.2024.3395423.

\bibitem{ref16} M.-T. Wu, H.-I. Lin, and C.-W. Tsai, ``A training-free neural architecture search algorithm based on search economics,'' \emph{IEEE Transactions on Evolutionary Computation}, vol. 28, no. 2, pp. 445--459, Apr. 2024, doi: 10.1109/TEVC.2023.3264533.

\bibitem{ref17} Y.-M. Kim, S. Song, B.-M. Koo, J. Son, Y. Lee, and J.-G. Baek, ``Enhancing long-term cloud workload forecasting framework: Anomaly handling and ensemble learning in multivariate time series,'' \emph{IEEE Transactions on Cloud Computing}, vol. 12, no. 2, pp. 789--799, Apr.--Jun. 2024, doi: 10.1109/TCC.2024.3400859.

\bibitem{ref18} F. Zhao, W. Lin, S. Lin, S. Tang, and K. Li, ``MSCNet: Multi-scale network with convolutions for long-term cloud workload prediction,'' \emph{IEEE Transactions on Services Computing}, vol. 18, no. 2, pp. 969--982, Mar.--Apr. 2025, doi: 10.1109/TSC.2025.3536313.

\bibitem{ref19} B. Feng and Z. Ding, ``Application-oriented cloud workload prediction: A survey and new perspectives,'' \emph{Tsinghua Science and Technology}, vol. 30, no. 1, pp. 34--54, Feb. 2025, doi: 10.26599/TST.2024.9010024.

\bibitem{ref29}
[dataset] Alibaba Group, ``Alibaba Cluster Trace Program: Cluster Trace
v2018,'' GitHub repository, 2018. [Online]. Available:
\url{https://github.com/alibaba/clusterdata/tree/master/cluster-trace-v2018}.
Accessed: Jul. 3, 2026.


\end{thebibliography}

\end{document}
